\chapter{Perception: the sector scans}
\label{ch:t-perception}

\code{lidar_sector_node} (navigation host, on Deep Robotics' own LiDAR driver there) turns every LiDAR
cloud into two sector scans: the
range of the nearest obstacle in each angular sector around the robot.

\begin{center}
\begin{tabularx}{\linewidth}{@{}>{\raggedright}p{0.27\linewidth}p{0.3\linewidth}X@{}}
  \toprule
  \textbf{Scan} & \textbf{Sectors} & \textbf{Read by} \\
  \midrule
  \topic{perception/safety_scan} & 72 over \ang{360}, \SIrange{0.1}{4}{m} & the protective field
    and the follower's guard \\
  \topic{perception/sector_scan} & 32 over \ang{180} ahead, \SIrange{0.1}{8}{m} & reactive
    avoidance \\
  \bottomrule
\end{tabularx}
\end{center}

\section{Which points are obstacles}

Each cloud goes through four filters before it is binned:

\begin{enumerate}
  \item \textbf{Noise}: Livox tags every point with its confidence. Adhesion points (bits 0--1,
    the glue of two surfaces) and particle points (bits 2--3: rain, dust) at the confidence of
    \param{noise_level} are dropped.
  \item \textbf{Self}: points inside the footprint grown by \param{self_margin} are the robot.
  \item \textbf{Levelling}: the cloud is rotated by the base's roll and pitch (from odom, TF), so
    heights count from the floor even while the body pitches in its gait. Without a recent
    roll and pitch (\param{tilt_max_age}) every sector is BLOCKED: an unlevelled cloud could put
    the floor in the band.
  \item \textbf{Height band}: only points between \param{height_min} (\SI{0.10}{m}) and
    \param{height_max} (\SI{1.60}{m}) above the floor are obstacles. With
    \param{ground_segmentation} the floor height is estimated per \param{ground_cell} first, for
    ramps and uneven yards.
\end{enumerate}

\subsection{Binning}

A sector's range is its $k$-th nearest point (\param{k} = 2): one stray return never blocks a
sector, and a sector needs $k$ points to count. Beyond \param{range_max} a sector is free; a
non-finite or blocked range is treated by every consumer as an obstacle at the robot.

\section{A blind LiDAR}

A covered window, a failed sensor or heavy rain empties the cloud. Fewer than \param{min_points}
(\num{100}) valid points in a cloud, or with \param{regions} and \param{min_region_points} fewer
than the minimum in any azimuth region, make the cloud blind: every sector BLOCKED, the field in
its BLIND state (Section~\ref{sec:t-field}). No cloud for \param{timeout_sec} does the same.

\begin{tipbox}[A partly covered sensor]
\param{min_region_points} catches what the total count misses: a bag on one side of the robot
leaves enough points overall. To find the normal count, set \param{min_region_points} very
high for a moment at the site: the warning \code{scans BLOCKED ... N points in region k} names the
weakest region's count. Then set the minimum at a third of it (a start-only parameter: restart).
\end{tipbox}

\section{The near-field blind zone}

A LiDAR sees nothing below its lowest beam. With the sensor at height $H$ above the floor and its
steepest beam $\delta$ below the horizontal, an object of height $h$ disappears from the cloud
once it is closer than
\begin{equation}
  r(h) = \frac{H - h}{\tan\delta}
\end{equation}
The robot walking toward a low obstacle sees it from afar, then loses it right where it matters.
The \textbf{obstacle memory} keeps every obstacle cell in the odom frame for \param{memory_sec}
(\SI{10}{s}) and adds it back to the scans while it is below the view, within \param{radius}.
The model takes $H$ from \param{base_height} and the LiDAR's mount (\code{lidar_pose}), and
$\delta$ from \param{lower_fov_deg} and the mount's tilt.

\begin{center}
\begin{small}
\begin{tabular}{@{}rrrr@{}}
  \toprule
  & \multicolumn{3}{c}{\textbf{blind radius} $r(h)$ for an object of height} \\
  \textbf{Lowest beam} $\delta$ & \SI{0.10}{m} & \SI{0.20}{m} & \SI{0.30}{m} \\
  \midrule
  \ang{7} (Mid-360, level) & \SI{3.1}{m} & \SI{2.3}{m} & \SI{1.5}{m} \\
  \ang{30} & \SI{0.66}{m} & \SI{0.48}{m} & \SI{0.31}{m} \\
  \bottomrule
\end{tabular}
\end{small}
\end{center}
\noindent ($H = \SI{0.48}{m}$, the sensor at base\_link height.) The X30's downward LiDARs see much
closer than a level sensor: set \param{lower_fov_deg} to the steepest beam of the cloud the node
reads, and check with a box of the height of \param{height_min} that the scan keeps it while the
robot walks up to it.

\begin{warnbox}[The memory forgets on a jump]
An odometry jump larger than the motion (\param{max_jump}) clears the memory: remembered cells
would land in the wrong place. A restarted odometry does the same.
\end{warnbox}

\section{Motion smear}

The Livox accumulates a cloud over $T = \SI{0.1}{s}$ while the robot moves. The cloud is read
as it arrives, so a point measured at the start of the sweep is off by up to
\begin{equation}
  \Delta \approx v\,T + \omega\,T\,r
\end{equation}
for a point at range $r$: \SI{7.5}{cm} straight ahead at \SI{0.75}{m/s}, \SI{14}{cm} at \SI{2}{m}
while walking at \SI{0.5}{m/s} and turning at \SI{0.45}{rad/s}. The protective field's
\param{margin} and the guard's \param{lateral_margin} must cover it. To remove it instead,
read FAST-LIO's undistorted cloud (\cmd{input.type: pointcloud2},
\cmd{input.pointcloud_topic: /cloud_registered_body}) with its frame in TF; it arrives one LIO
step later.

\section{Tuning the band}

\begin{description}
  \item[\param{height_min}] above the floor's unevenness and the gravel, below the lowest thing that
    can trip the robot. Points of a slope rise into the band: enable
    \param{ground_segmentation} on ramps rather than raise the bottom.
  \item[\param{height_max}] above the highest part of the robot with its payload (the PTZ camera on
    its mast): anything lower that hangs over the route must be seen.
  \item[\param{range_max} (safety)] at least the field's distance at full speed plus the footprint:
    \SI{1.62}{m} + \SI{0.50}{m} for the X30; \SI{4}{m} leaves room.
  \item[\param{num_sectors}] 72 gives \ang{5} sectors: at \SI{2}{m} a sector is \SI{17}{cm} wide.
    More sectors resolve a narrow gap better and need more points each.
\end{description}
